% Starter deck for the project presentation. Shares macros with the paper.
\documentclass[11pt,aspectratio=169]{beamer}
\usetheme{metropolis}
\usepackage[T1]{fontenc}
\usepackage{amsmath, amssymb}
\usepackage{booktabs}
\input{../paper/macros}

\title{Does Random Mean Secure?}
\subtitle{Randomness, Entropy and Perfect Secrecy}
\author{tdvnm}
\date{\today}

\begin{document}

\maketitle

\begin{frame}{Two meanings of ``random''}
  \begin{enumerate}
    \item A sequence \emph{looks} random: balanced, right run lengths,
          nothing repeats. \hfill (statistics)
    \item A sequence \emph{is} unpredictable: knowing everything else, you
          still cannot guess the next bit. \hfill (cryptography)
  \end{enumerate}
  \pause
  \begin{block}{The question}
    If the randomness behind a key gets worse, how much more does the
    ciphertext reveal about the message?
  \end{block}
\end{frame}

\begin{frame}{Perfect secrecy}
  Shannon (1949): a cipher is perfectly secret if for every $m$ and $c$
  \[
    \Prob{\msg = m \mid \ctxt = c} = \Prob{\msg = m}.
  \]
  Equivalently, $\MI{\msg}{\ctxt} = \Ent{\msg} - \CEnt{\msg}{\ctxt} = 0$.
  \medskip

  For $\ctxt = \msg \xor \key$ with $\key$ uniform on $\bits^n$ and
  independent of $\msg$:
  \[
    \Prob{\ctxt = c \mid \msg = m} = \Prob{\key = m \xor c} = 2^{-n},
  \]
  which does not depend on $m$, so Bayes' rule gives perfect secrecy.
\end{frame}

\begin{frame}{The main experiment}
  Fix the cipher $\ctxt = \msg \xor \key$, $\msg \in \bits^2$ lopsided.
  Only change the key source.
  \begin{center}
    \begin{tabular}{@{}lll@{}}
      \toprule
      Key & Expected $\MI{\msg}{\ctxt}$ \\
      \midrule
      A. uniform         & $\to 0$ \\
      B. one value w.p.\ $0.7$ & $> 0$ \\
      C. bits $1$ w.p.\ $q \in [0.5, 0.95]$ & grows with $q$ \\
      \bottomrule
    \end{tabular}
  \end{center}
  \medskip
  Main figure: leakage against key bias, simulation vs.\ exact.
\end{frame}

\begin{frame}{Takeaway}
  \begin{center}
    \Large Good randomness is necessary for secrecy.\\
    But looking random is not enough to make something secure.
  \end{center}
\end{frame}

\end{document}
